% !TEX root = ./pkArticle.tex
\mode*

\subsection{Elimination, Context sensitive Halbwertszeit}

Mit der {\it Halbwertszeit} (\HWZ) wird die Geschwindigkeit von (biologischen) Prozessen beschrieben. Beim radiokativen Zerfall ist die Aktivit�t nach einer \HWZ  
halbiert. Nach zwei \HWZ betr�gt die Aktivit�t noch die H�lfte der H�lfte d.h. ein Viertel.

Diese �berlegungen lassen sich nicht direkt auf den Konzentrationsabfall von Medikamenten �bertragen. Die einzelnen Umvertelungsprozesse k�nnen ebenfalls mit \HWZ beschrieben werden doch ist die Bedeutung nicht dieselbe wie beim radiokativen Zerfall. Deshalb soll die {\it Elimination} etwas allgemeiner betrachtet werden.

\begin{frame}
\frametitle{Elimination}
\begin{itemize}[<+->]
\item
Elimination aus K\"orper oder Elimination aus Blut?
\item
Konzentrationsabfall: Umverteilung und Metabolisierung
\item
Es kann aktive Metaboliten geben
\item
Fettl\"osliche Medikamente werden in Leber zu inaktiven, wasserl\"oslichen Metaboliten umgewandelt
\item
Ausscheidung der wasserl\"oslichen Substanzen durch Niere
\item
Pharmakokinetik {\bf{quantifiziert}} die Geschwindikeit der Elimination
\end{itemize}

\note<1>{

\begin{itemize}
\item
Elimination wird meist als \enquote{Kleinerwerden der Konzentration} des Medikamentes im Blut quantifiziert!
\item
Meist denkt man dabei an die Elimination des Medikamentes aus dem \enquote{Wirkort}, d.h. \enquote{Kleinerwerden} der Wirkung.
\item
Elimination ist aber im Grunde genommen das \enquote{Verschwinden} des Medikamentes aus dem K�rper
\end{itemize}

}
\end{frame}



\frame{
\mode<presentation>{
\begin{tikzpicture}[remember picture, overlay]
\node[inner sep=0pt,xshift=0cm] at (current page.center){
\tikz \draw[step=2mm,black!50] (0,0) grid (126mm,94mm);
};
{
\node[coordinate] at (1cm,-3.55cm) (start) {};
\node[coordinate] at (1cm,3cm) (end) {};
\node[coordinate] at (10cm,-3.55cm) (endy) {};
\draw[->,thick] (start)--(end);
\draw[->,thick] (start)--(endy);
};
\end{tikzpicture}
}
\frametitle{Halbwertszeit}
\note<1>{
\begin{itemize}
\item
Halbwertszeit wird oft als \enquote{Mass} f�r Eliminationsgeschwindigkeit verwendet
\item
Zeigen was Halbwertszeit bedeutet.
\item
Zeichnen eines Kompartimentes: Elimination ver�ndert sich. (Es wird immer die H�lfte eliminiert)
\item
Wesentlich: $t_{\frac{1}{2}}$ bleibt konstant
\item
Mathematische Beziehung zwischen Clearance und $t_{\frac{1}{2}}$
\end{itemize}
}
\infina
}


\mode<article>{
\vspace{1cm}
\begin{tikzpicture}[auto]
\node[inner sep=0pt,xshift=0cm] at (current page.center){
\tikz \draw[step=2mm,black!50] (0,0) grid (140mm,70mm);
};
\end{tikzpicture}
}


\mode<article>{\pagebreak}


\begin{frame}
\frametitle{Halbwertszeit: Mehrkompartiment Modelle}
\note<1>{
\begin{itemize}
\item
Zeichnen: Drei Kompartiment Modell - mit drei Clearance Werten
\item
Gleichzeitig wirken drei Clearance. Zu Beginn $\Delta Conc$ gr\"osser
\item
Es gibt definierte $t_{\frac{1}{2}}$: Achtung $2 * t_{\frac{1}{2}} \neq \frac{C_{0}}{4}$ !
\item
Wenn Zentrales Kompartiment bei 1 Komp aufgef\"ullt wird, immer wieder dieselbe Ausgangssituation\\
bei Mehrkompartiment ist das ganz anders. Die Halbwertszeit wird anders sein.
\item
Mit Cylinder zeigen. 1 Komp auff\"ullen, 3 Komp auff\"ullen.\\
10 l, 20l, 200l, 2,1,0.5 :  Bolus: von 50 = Konz 5; Nach t1/2 wieder Bolus 25 ...
\end{itemize}
}
\infina
\end{frame}


\mode<article>{
\vspace{1cm}
\begin{tikzpicture}[auto]
\node[inner sep=0pt,xshift=0cm] at (current page.center){
\tikz \draw[step=2mm,black!50] (0,0) grid (140mm,70mm);
};
\end{tikzpicture}
}


\begin{frame}
\frametitle{Halbwertszeit bei Mehrkompartiment Kinetik}
\begin{itemize}
    \item Gleichzeitg verschieden schnelle Umverteilung
	\begin{itemize}
    	\item Schnelle Umverteilung in Muskel (gut durchblutete Gewebe)
    	\item Langsamere Umverteilung in Fettgewebe
    	\item Elimination (Leber, Niere)
	\end{itemize}
    \item Die Prozesse haben {\it eigene} Halbwertszeit
\end{itemize}
\note<1>{.}
\end{frame}


Die Bedeutung dieser Halbwertszeiten in Bezug auf den Konzentrationsabfall �ndert sich mit der Dauer der Medikamentenzufuhr. Wenn nach mehreren Bolusdosen die peripheren Kompartimente \enquote{aufgef�llt} sind, ist der initiale Konzentrationsabfall kleiner als nach der ersten Dosis.



\begin{frame}
\frametitle{Halbwertszeit ist von Dauer der Infusion abh\"angig! (Kontext sensitive Halbwertszeit)}
\begin{itemize}
\item
Dauer der Infusion ist der \enquote{Kontext}
\item
Je l�nger die Infusion dauert, desto mehr Medikament ist im K�rper umverteilt, umso l�nger dauert die Elimination.
\item
Die aktuell g�ltige Halbwertszeit kann berechnet werden (Mit PK Modell)
\end{itemize}
\note<1>{
\begin{itemize}
\item
Excel: Dauer der Infusion: Halbwertszeit wird l\"anger. SimInfCSHT.xls
\item
Excel sheet mit Kurven: CSDTSimPropRemi.xls
\end{itemize}
}
\infina
\end{frame}


Praktisch bedeutet dies, dass je l�nger eine An�sthesie dauert, desto fr�her (vor Ende der OP) muss die Medikamentenzufuhr (Propofol) gestoppt werden, damit die Patientin rechtzeitig erwacht.


